\begin{lemma}[Scheide replacement-pair switch invariance (8)]
For replacement data of a defective color $\delta$, the switched coloring
$\varphi_\delta$ agrees with $\varphi$ on every edge identity incident with
$U_\delta=\{u^1_\delta,u^2_\delta\}$.
\end{lemma}
\begin{proof}
Suppose that $u^i_\delta$ lay on the switched
$(\gamma(\delta),\delta)$ component.  Its unique incident
$\delta$-edge $f^i_\delta$ would then be an edge-node of that component,
and its other endpoint $v^i_\delta\in V(T)$ would lie on the same component.
This tree vertex differs from $v^0_\delta$, since the three boundary edge
identities in \noderef{ScheideReplacementData} are distinct and a proper
coloring has at most one incident $\delta$-edge at $v^0_\delta$.  This
contradicts \noderef{ScheideDefectiveTailExternal}, which says that
$v^0_\delta$ is the component's only tree vertex.  Thus neither replacement
vertex-node belongs to the component.

If an edge identity incident with one of these vertices belonged to the
component, incidence adjacency would put the vertex-node in the same
component.  Hence every such edge-node is outside the component and is
fixed by the exact definition \noderef{TwoColorComponentSwap}.  This is
the asserted equality for every incident edge identity.
\end{proof}
